\section{Policy extensions}\label{si:policy}

\paragraph{Half-Earth and reservation.} The Half-Earth proposal \citep{Wilson2016} caps human-usable annual biocapacity account value (gha) at half of the total. In the two-land setting this is a reservation $E\le\sigma B$ with dimensionless $\sigma$ the human-available share, enforced either as a cap on the fast book or on the carrying capacity; with $\sigma_f,\sigma_c<1$ in Eq.~(4) the shortfall is computed against the reserved flows, so conversion begins earlier for a given demand and the ceiling of Proposition~2 is lowered to $\sigma$-weighted values. On the nominal path a reservation that keeps $E\le\sigma B$ succeeds with no liquidation. This is a nominal-path consequence until it is certified against a declared disturbance class (persistent productivity, $b_f$ or $\rho_c$ shocks) under a preregistered retention rule (a governance module is retained only if it improves a declared protection-and-supply score), with the model's declared fit defects converted into an erosion margin on the closed loop's contraction rate; until then it is a property of the nominal calibration, not a robustness claim.

\paragraph{Typed floors and freeze conversion.} Reserve $A_r\ge A_r^{\min}$, capital viability $A_c\ge\Acmin$, arable ceiling $A_f\le A_f^{\max}$ and $K\ge K_{\min}$ are separate constraints. Freezing conversion (capping $A_f$ at $A_f^{\max}$ or flooring $A_c$ at $\Acmin$) is the single operator that guarantees $A_c\ge\Acmin$; the conversion rate saturates as $A_c\to\Acmin$ by construction of Eq.~(5). A policy that observes only the aggregate $B$ cannot guarantee capital-land integrity: the failure is informational (aggregation relative to the typed floor), not dynamical.

\paragraph{Institutional signature.} In the mobilizing/protective taxonomy of the delay companion \citep{Abaee2026a}, freeze-conversion and surplus-gated restoration are protective; the mobilizing sign (demand responding to a perceived shortfall by converting more) is the destabilizing one; deficit-gated restoration is not unambiguously protective because it can deepen the short-run shortfall (Section~\ref{si:gates}).

\paragraph{Weak and strong sustainability reading.} In the material-cycle reading of the ledger companion \citep{Abaee2026c}, weak sustainability is the idealized regime in which substitution and regeneration jointly redistribute matter so that it is used as it arises -- substitution admissible only through an identified physical pathway, never merely because a weighted aggregate permits it -- and strong sustainability is ``the regime in which that closure fails -- either because no identified physical pathway re-routes the extracted matter, or because depletion outruns the rate at which such a pathway could be deployed'', so that a separately binding floor on a critical stock binds. The composite annual account value $B$ (gha) is the weak-sustainability index: the linear, perfectly substitutable member of the aggregator family of the assessment companion \citep[Lemma~A of][]{Abaee2026b}; it can satisfy its aggregate floor while the typed floor on $A_c$ is violated by exactly the amount that the fast-land yield gains against the capital loss, and in that family no positive-elasticity aggregator is uniformly safe (its Theorem~S2(ii)). Under exact-tube semantics the intermediate passage during which $A_c$ falls is itself a violation even if the endpoint recovers.

\paragraph{Reporting rule.} A growth headline for total biocapacity carries no information about whether cropland growth is compensating for stagnation or offsetting changes in the other components. The informative assessment states, at every level at which an aggregate is used: component endpoint ratios, log contributions, attribution shares, data vintage and endpoint sensitivity.

\section{The one-stock comparator}\label{si:onestock}

Earlier versions of this work developed a single-stock model that the two-land model supersedes. It is retained here as a comparator so that quantities that circulate for it are not mistaken for two-land results. The comparator has three states ($A$, $P$, $D$) and one algebraic relation:
\begin{align}
G(A)&=\rho A(1-A/A_{\max}),\qquad B=bA+b_GG(A),\qquad E=eP,\\
\dot A&=G\big(A(t-\tau_g)\big)-\pos{E-\sigma bA}/b_G,\\
\dot P&=rP\big(1-P/K(t-\tau_p)\big),\qquad K=B/e,\\
\dot D&=\pos{E-B}-\eta D,\qquad b=(b_0+T_b)e^{-\alpha D},
\end{align}
Here $A$ and $A_{\max}$ are comparator-specific stock units and $G(A)$ has units stock unit\,yr$^{-1}$; consequently $b,b_0,T_b,\Delta b$ are gha per stock unit and $b_G$ is gha\,yr per stock unit, so $bA$ and $b_GG(A)$ are each gha. $B$ and $E$ are annual GFN/NFA account amounts in gha, $P$ and $K$ are in cap, and $e$ is gha\,cap$^{-1}$. The shares $\sigma$ are dimensionless. Time and the delays are in yr; $\rho,r,\eta$ and the one-stock logistic-wave steepness $\kappa$ are rates in yr$^{-1}$ ($\kappa$ is distinct from the two-land $\kappa_w$); $D$ is a leaky time-integral in gha\,yr, so $\dot D$ and $\eta D$ are gha; and $\alpha$ is gha$^{-1}$\,yr$^{-1}$. The figures share a numerical ordinate only for illustration; $A$, $P$, $B$ and $D$ retain distinct units.

With demand met first from the flow $bA$, the shortfall liquidates the standing stock at rate $1/b_G$. Its properties, each of which was re-derived under the two-land structure and found not to carry over, are the following.
\begin{itemize}[leftmargin=1.5em]
\item \emph{MSY and fold.} $B(A)$ has an interior maximum (the MSY of Schaefer, 1954) only in the capital-dominated regime $b_G\rho>b$; a fixed liability above $B_{\max}$ produces a saddle-node, the catastrophic-shift structure of \citet{Scheffer2001}. At the comparator baseline ($b_G\rho=0.04<b=0.5$ in gha per comparator-specific stock unit) the fold is absent. In the two-land model $Y_c$ is linear in $A_c$ and $B_{eq}(A_c)$ is monotone; there is no fold in any regime (Proposition~2).
\item \emph{The equilibrium family $P=B(A)/e$ and $\det J=0\,\mathrm{yr}^{-3}$.} Under full harvest ($\sigma=1$) the comparator has a neutral continuum of equilibria on which the Jacobian is exactly singular. The two-land model has a one-parameter family in the \emph{land allocation} (indexed by $A_c$ on the no-conversion face, with $A_r/A_f=\eta_f/\mu$), not a free $P$, and no identically singular Jacobian.
\item \emph{The $+0.6$ eigenvalue and the two-gain Hopf classification.} The comparator's deficit-region linearization has a positive real root ($a_{11}=G'(A^*)+b/b_G\approx0.608$--$0.625$\,yr$^{-1}$ at its baseline) and a dimensionless comparator-specific two-gain rule $\chi_{\mathrm{1s}}=q/(\rho-2q)$ (here $q$ is the comparator's feedback-rate coefficient in yr$^{-1}$, not the two-land quality index) selecting which delay destabilizes ($\tau_g$-only for $\chi_{\mathrm{1s}}>1$, $\tau_p$-only for $\chi_{\mathrm{1s}}<1$, with critical delays $\tau_m\approx85.4$\,yr and $\tau_p\approx231$\,yr). None of this is a two-land quantity: the two-land deficit-side spectrum is real and negative ($-0.054\,\mathrm{yr}^{-1}$, $-1.0\,\mathrm{yr}^{-1}$) and constant along the branch.
\item \emph{The regeneration cliff and the masking window.} The comparator shows a recover/collapse boundary at $\tau_g\approx18$--$20$\,yr and a productivity-illusion window of $\approx5.4$\,yr ($A:1.00\to0.86$ comparator-specific stock units while $B:0.578\to0.647$\,gha; the mask vanishes for initial deficits $\gtrsim0.075$\,gha). In the two-land model the boundary is a demand threshold insensitive to $\tau_g$ (Table~\ref{tab:si_delay}) and the composition window is set by $\lambda_1$ and $\dconv$ (paper, Section~3.1). The dimensionless recovery fractions $0.54\to0.52\to0.50$ for $\tau_g=5\to25\to40$\,yr quoted in earlier versions were computed with an area-based intermediate formulation and are superseded by the identical-to-$10^{-3}$ result of Table~\ref{tab:si_delay}.
\end{itemize}
The four selected comparator figures reproduced here are numbered \ref{fig:si_s5}--\ref{fig:si_s12}. The broader historical supplementary figure set (S1--S15) remains preserved in the repository at \texttt{agent 2 productivity illusion/supplementary/FIGURES/}; only figures used by the current SI are included in this package.

\begin{figure}[htbp]
\centering
\includegraphics[width=0.8\textwidth]{S5_demo_overshoot_run.png}
\caption{One-stock comparator: top panel shows stock $A$ (model-specific stock units), population $P$ (cap) and annual biocapacity $B$ (gha); the lower panel shows leaky debt $D$ (gha\,yr). The footprint $E=eP$ is a driver and is not plotted. The shared numerical ordinate is illustrative, not a common physical scale. Comparator only; not a two-land result.}
\label{fig:si_s5}
\end{figure}

\begin{figure}[htbp]
\centering
\includegraphics[width=0.8\textwidth]{S6_masking_window.png}
\caption{One-stock comparator: the productivity-illusion window ($A$ falls from $1.00$ to $0.86$ model units while $B$ rises from $0.578$ to $0.647$ gha; window $\approx5.4$\,yr; the mask vanishes for deficits $\gtrsim0.075$ gha). The two-land composition window of the paper (Fig.~1a) is a different object with a different mechanism.}
\label{fig:si_s6}
\end{figure}

\begin{figure}[htbp]
\centering
\includegraphics[width=0.8\textwidth]{S7_recovery_vs_collapse.png}
\caption{One-stock comparator: recovery versus collapse as a function of regeneration delay $\tau_g$ (yr; the $\approx18$--$20$\,yr cliff). The plotted series use a shared numerical ordinate only for illustration: stock $A$ is in comparator-specific units, biocapacity $B$ in gha and population $P$ in cap. Not present in the two-land model.}
\label{fig:si_s7}
\end{figure}

\begin{figure}[htbp]
\centering
\includegraphics[width=0.8\textwidth]{S12_basin_baseline.png}
\caption{One-stock comparator: basin of recovery at the comparator baseline. The axes are initial population $P_0$ (cap) and initial stock $A_0$ (comparator-specific stock units); the equilibrium guide $P=B(A)/e$ is in cap, with annual account value $B$ in gha and $e$ in gha\,cap$^{-1}$. In the two-land model the corresponding object is the continuous drawdown-depth surface of Fig.~1c and Table~\ref{tab:si_drawdown}.}
\label{fig:si_s12}
\end{figure}

\section{Controlled technology-wave experiment}\label{si:tbwave}

This controlled extension tests whether an active $T_b$ wave can produce a larger leaky-debt peak than an otherwise matched no-wave run. The pair uses endogenous demand $E=eP$ (so the surplus-gated $R_{rc}$ is active only when $B>E$), $A_{tot}=2.5$\,ha, $A_{f0}=0.6$\,ha, $A_{c0}=0.9$\,ha, $q_0=0.9$, $D_0=0.0$\,gha\,yr, and $P_0=1.0$ or $1.2$\,cap. Shared parameters are $b_{f0}=0.85$\,gha\,ha$^{-1}$, $b_c=0.05$\,gha\,ha$^{-1}$, $b_{G,c}=0.8$\,gha\,yr\,ha$^{-1}$, $\rho_c=0.08$\,yr$^{-1}$, $\alpha=0.03$\,gha$^{-1}$\,yr$^{-1}$, $\eta=0.05$\,yr$^{-1}$, $\eta_f=0.05$\,yr$^{-1}$, $\mu=0.06$\,yr$^{-1}$, $e=0.55$\,gha\,cap$^{-1}$, $r=0.02$\,yr$^{-1}$, $\tau_g=20$\,yr, $\tau_p=25$\,yr, $\tau_{conv}=1$\,yr, and $\chi_r=0.10$\,gha$^{-1}$\,yr$^{-1}$. The wave parameters are $\Delta b=1.0$\,gha\,ha$^{-1}$, $\kappa_w=0.1$\,yr$^{-1}$, and $t_{wave}=40$\,yr; the horizon is $T=250$\,yr. Within each pair, the physical initial state, parameters, wave-free pre-run history, and solver are identical; only the $T_b$ switch differs. The logistic is already nonzero at $t=0$ ($T_b(0)=0.017986$\,gha\,ha$^{-1}$), so the derived initial account is $B(0)=0.560184$\,gha without the wave and $0.570976$\,gha with it. This small initial difference is part of the declared forcing, not a change to $A_f,A_c,q,P$, or $D$; $T_b$ reaches half its increment at the marked midpoint $t_{wave}=40$\,yr.

Figure~\ref{fig:si_wave_debt} shows the $P_0=1.0$ pair. Its all-time debt maximum occurs after the wave midpoint in both runs: $D_{max}=0.008126$\,gha\,yr at $86.15$\,yr without the wave, versus $0.020205$\,gha\,yr at $115.60$\,yr with it. At $P_0=1.2$, by contrast, the all-time maximum occurs early and is lower with the wave, while the maximum over $t\ge t_{wave}$ is higher. Table~\ref{tab:si_wave_debt} keeps these all-time and post-midpoint definitions separate. We call $t\ge t_{wave}$ the post-midpoint window; it is not an assertion that a logistic wave has ended at $t_{wave}$. The full metrics and trajectories are deposited as \texttt{work/tb\_wave\_debt\_metrics.csv}, \texttt{work/tb\_wave\_debt\_timeseries\_P0\_1p0.csv}, and \texttt{work/tb\_wave\_debt\_timeseries\_P0\_1p2.csv}; \texttt{work/tb\_wave\_debt\_scan.py} regenerates them and the figure.

\begin{table}[htbp]
\centering\small
\caption{Matched technology-wave/no-wave debt maxima. Entries are $D_{max}$ in gha\,yr, with the time of the maximum in years in parentheses; the ratio is wave/no-wave. ``All'' is the maximum over $0\le t\le250$\,yr; $t\ge40$ is the post-midpoint maximum. The $P_0=1.0$ all-time peaks both occur after the midpoint.}
\label{tab:si_wave_debt}
\begin{tabularx}{\textwidth}{r l r l >{\centering\arraybackslash}X >{\centering\arraybackslash}X r}
\toprule
$P_0$ (cap) & quality update & $\Delta t$ (yr) & window & no wave $D_{max}$ (time) & wave $D_{max}$ (time) & ratio \\
\midrule
$1.0$ & delayed & $0.050$ & all & $0.008126\ (86.15)$ & $0.020205\ (115.60)$ & $2.487$ \\
$1.0$ & delayed & $0.025$ & all & $0.008105\ (86.20)$ & $0.020195\ (115.65)$ & $2.492$ \\
$1.0$ & current & $0.050$ & all & $0.008326\ (85.70)$ & $0.020168\ (115.60)$ & $2.422$ \\
$1.2$ & delayed & $0.050$ & all & $0.061218\ (1.65)$ & $0.051037\ (1.45)$ & $0.834$ \\
$1.2$ & delayed & $0.050$ & $t\ge40$ & $0.010854\ (81.70)$ & $0.025184\ (112.55)$ & $2.320$ \\
$1.2$ & delayed & $0.025$ & all & $0.060732\ (1.625)$ & $0.050582\ (1.475)$ & $0.833$ \\
$1.2$ & delayed & $0.025$ & $t\ge40$ & $0.010808\ (81.775)$ & $0.025165\ (112.625)$ & $2.328$ \\
\bottomrule
\end{tabularx}
\end{table}

The $P_0=1.0$ result is a declared sensitivity case, not a refit or a universal effect: wave/no-wave all-time peak ratios persist under step halving ($2.487$ at $\Delta t=0.05$\,yr; $2.492$ at $0.025$\,yr) and under the current-quality solver ($2.422$ at $0.05$\,yr). At the reference $P_0=1.2$, the corresponding all-time ratios are $0.834$ and $0.833$, whereas the post-midpoint ratios are $2.320$ and $2.328$. All reported debt maxima change by less than $0.9\%$ under step halving. Maximum area-conservation residual is $4.44\times10^{-16}$\,ha; the minimum reserve in the $P_0=1.0$ wave run is $0.286$\,ha, above the $0.10$\,ha floor. In this sensitivity pair $A_c$ does not fall below its initial $0.9$\,ha, so the experiment isolates a debt rebound and is not evidence of the capital drawdown shown in the main frozen-demand run. The regression test also confirms that an integer input $D_0=0$ still produces a floating-point debt state and finite debt accumulation. $D$ is the model's leaky time-integral of shortfall, not an observed NFA or physical-stock variable. These conditional trajectories demonstrate a possible delayed rebound, while the $P_0=1.2$ contrast shows why the all-time peak cannot be substituted for the post-midpoint response.

\begin{figure}[!htbp]
\centering
\includegraphics[width=0.92\textwidth]{fig_tb_wave_debt}
\caption{Matched endogenous-demand technology-wave/no-wave simulation at $P_0=1.0$\,cap. Parameters and initial conditions are listed in SI~\S S8; only $T_b$ is switched on or off. (a) The active logistic increment $T_b(t)$, with the dotted line at its midpoint $t_{wave}=40$\,yr; (b) annual account values $B$ and $E$ (gha); (c) population $P$ (cap); (d) leaky model debt $D$ (gha\,yr). Markers denote all-time maxima, which occur at $86.15$\,yr without the wave and $115.60$\,yr with it; peak debt is $0.008126$ versus $0.020205$\,gha\,yr. The plotted wave begins with its small logistic tail at $t=0$ and approaches $\Delta b=1.0$\,gha\,ha$^{-1}$; the dotted line marks the midpoint, not an end date.}
\label{fig:si_wave_debt}
\end{figure}

\section*{Supplementary references}
